
% SGA 3，报告 VII_B；英文重建本第一轮的中文译本。
% 范围： 各节 4.2 与 4.2.1 全文 依据文献局部页 75,
% 止于下节之前： 节 4.2.2.
% 合并阅读本页： 588. 重版印刷页： 577.
% 范围内原版边码： 543.
% 依据： Polo--Gille Exp7B-13oct24.pdf,
% SHA-256 D303690E0195654BBE7C3DAA5D89EA23E24CD709A3575217234AC19AE8836694.
% 参照了 Jacob Reinhold 的 CC BY 4.0 Markdown，作为对照来源。
% 本组件无图；所有数学内容均为原生 LaTeX。
\subsection*{4.2}
\phantomsection
\label{sga3:VIIB:4.2}
显然，形式 \(k\)-群 \(G\) 的 Lie 代数 \(\operatorname{Lie}(G)\) 是 \(\mathbf H(G)\) 的 \(p\)-Lie 子代数\textup{（参见\ \SGARefLink{sga3:VIIA:2.3}{2.3}）}。事实上，可立即归约到 \(k\) 是 Artin 环的情形。此时，\(\operatorname{Lie}(G)\) 是 \(\mathbf H(G)\) 中满足
\[
  \varphi_G(x\otimes1+1\otimes x)=\Delta_G(x),
\]
的元素 \(x\) 构成的子集，这里采用 \SGARefLink{sga3:VIIB:2.3}{2.3} 和 \SGARefLink{sga3:VIIB:2.6}{2.6}\textup{(c)} 的记号。故
\[
\begin{aligned}
  \varphi_G(x^p\otimes1+1\otimes x^p)
  &=
  \varphi_G\bigl((x\otimes1+1\otimes x)^p\bigr)
  \\
  &=
  \varphi_G(x\otimes1+1\otimes x)^p
  \\
  &=
  \Delta_G(x)^p
  =
  \Delta_G(x^p).
\end{aligned}
\]
\subsubsection*{4.2.1}
\phantomsection
\label{sga3:VIIB:4.2.1}
反过来，每个 \(\mathbf O_k\) 上的 \(p\)-Lie 代数 \(\mathbf L\) 都定义一个 \(k\)-群函子。用 \(U_p(\mathbf L)\) 表示将 \(\operatorname{Alf}/k\) 的每个对象 \(C\) 映为 \(C\) 上 \(p\)-Lie 代数 \(\mathbf L(C)\) 的限制包络代数
\[
  U_p\bigl(\mathbf L(C)\bigr)
\]
的函子\textup{（VII\textsubscript{A}，\SGARefLink{sga3:VIIA:5.3}{5.3}）}。由 VII\textsubscript{A}，\SGARefLink{sga3:vii:corollary:5-4}{5.4}，\(U_p(\mathbf L)\) 是 \(\mathbf O_k\) 上的双代数，所以由 \SGARefLink{sga3:VIIA:2.2}{2.2}，它确定 \(k\)-群函子
\[
  \operatorname{Spf}^{*}U_p(\mathbf L),
\]
此后将其记为 \(\mathscr G_p(\mathbf L)\)。
于是，\(\mathscr G_p(\mathbf L)(C)\) 是增广为 \(1\) 且满足下式的元素
\[
  z\in U_p\bigl(\mathbf L(C)\bigr)
\]
所构成的群：
\[
  \Delta_{U_p(\mathbf L(C))}(z)=z\otimes z.
\]

